\section{Paths, compactness, and local sections for groups of matrices}

This section proves that the invertible, the unitary, and the real orthogonal
matrices of a fixed determinant are path connected, that the matrices unitary
between two coordinate spaces form a compact set, and that a path of unitaries
that is continuous modulo local unitaries $x\otimes y$ lifts to a continuous
path. The last statement replaces the slice theorem of Montgomery and Yang in
the proof of \cite[Corollary~IV.7]{Cirac2017MPU}; the construction is
described in \cite{gap:mpu_standard_form_path_cross_section}.

\begin{lemma}[Paths of invertible matrices]
    \label{lem:invertible_matrices_joined}
    \lean{Matrix.joinedIn_isUnit_det_one}
    \leanok
    Every invertible complex square matrix $X$ is joined to the identity by a
    continuous path of invertible matrices.
\end{lemma}

\begin{proof}
    \leanok
    For $z\neq0$, the matrix $(1-z)I+zX=-z\bigl(-z^{-1}I-(X-I)\bigr)$ is
    singular exactly when $-z^{-1}$ is an eigenvalue of $X-I$. Hence the
    line $z\mapsto(1-z)I+zX$ is singular at only finitely many points, none of
    them $0$ or $1$. The complement of a countable subset of $\mathbb C$ is
    path connected, so a path from $0$ to $1$ avoiding these points gives the
    required path of invertible matrices.
\end{proof}

\begin{lemma}[Paths of unitary matrices]
    \label{lem:unitary_matrices_joined}
    \lean{Matrix.joined_one_unitaryGroup, Matrix.IsUnitaryBetween.joinedIn}
    \leanok
    Every unitary complex square matrix is joined to the identity by a
    continuous path of unitary matrices. Consequently, any two linear maps
    that are unitary between the same pair of finite-dimensional coordinate
    spaces are joined by a continuous path of such maps.
\end{lemma}

\begin{proof}
    \leanok
    Let $u$ be unitary. Its spectrum is a finite subset of the unit circle,
    so there is $z$ on the unit circle with $-z$ not an eigenvalue of $u$.
    The spectrum of $z^{-1}u$ then avoids $-1$, hence
    \begin{align}
        \|z^{-1}u-1\| & <2,
    \end{align}
    and $z^{-1}u$ is joined to the identity by the exponential of a
    multiple of its self-adjoint logarithm. The scalar $z$ is joined to $1$
    along the unit circle, and the product of the two paths joins $1$ to
    $u$. For maps $a,b$ unitary between the same spaces, $ba^\dagger$ is
    unitary; a path $g(t)$ from $1$ to $ba^\dagger$ gives the path $g(t)a$
    from $a$ to $b$.
\end{proof}

\begin{lemma}[Paths of orthogonal matrices]
    \label{lem:mpu_orthogonal_gate_interpolation}
    \lean{Matrix.IsUnitaryBetween.joinedIn_real_of_det_eq,
        Matrix.joined_unitaryGroup_real_of_det_eq,
        Unitary.joined_one_linearIsometryEquiv_symm}
    \leanok
    Two real orthogonal matrices with the same determinant are joined by a
    continuous path of real orthogonal matrices. Consequently, two complex
    matrices with real entries that are unitary between the same coordinate
    spaces, and whose determinants agree when read through one
    identification of those spaces, are joined by a path of such matrices.
    This is the interpolation step of the proof of
    \cite[Theorem~\texttt{conjphs}]{Cirac2017MPU}, of which only path
    connectedness is used.
    % Source lines: paper_v2.tex 1438--1442.
\end{lemma}

\begin{proof}
    \leanok
    By the Cartan--Dieudonn\'e theorem an orthogonal map is a product of
    reflections in hyperplanes, an even number of them when its determinant
    is one. In dimension at least two, any two such reflections are joined by
    moving the unit normal along the sphere, so the product is joined to an
    even power of one reflection, which is the identity.
\end{proof}

\begin{lemma}[Compactness of the unitaries between two coordinate spaces]
    \label{lem:unitary_between_compact}
    \lean{Matrix.isCompact_setOf_isUnitaryBetween}
    \leanok
    For finite sets $M$ and $N$, the matrices $A\in\C^{M\times N}$ with
    $A^\dagger A=I_N$ and $AA^\dagger=I_M$ form a compact set.
\end{lemma}

\begin{proof}
    \leanok
    Every entry of such a matrix has modulus at most one, so the set lies in
    a product of closed discs, which is compact. Both equations are closed
    conditions, since multiplication and the adjoint are continuous.
\end{proof}

\begin{lemma}[Nearest Kronecker product]
    \label{lem:nearest_kronecker_product}
    \lean{Matrix.kroneckerProjector, Matrix.nearestKronecker,
        Matrix.kroneckerGapSet, Matrix.continuous_nearestKronecker,
        Matrix.isOpen_kroneckerGapSet, Matrix.nearestKronecker_kronecker_mul,
        Matrix.exists_eq_kronecker_nearestKronecker,
        Matrix.nearestKronecker_kronecker_of_unitary,
        Matrix.kronecker_mem_kroneckerGapSet}
    \leanok
    \uses{lem:matrix_realignment_kronecker}
    Let $0<a<b<\ell r$, and let $f$ be the continuous function equal to $0$ on
    $(-\infty,a]$, to $1$ on $[b,\infty)$, and affine in between. For
    $M\in\mathrm M_{\ell r}$ put $P(M)=f\bigl(R(M)R(M)^\dagger\bigr)$, where
    $R$ is the realignment, and let $N(M)$ be the matrix with
    $R(N(M))=P(M)R(M)$. Let $W$ be the set of $M$ such that the spectrum of
    $R(M)R(M)^\dagger$ does not meet $[a,b]$ and
    $\operatorname{tr}P(M)\in(1/2,3/2)$. Then:
    \begin{enumerate}
        \item $N$ is continuous and $W$ is open;
        \item $N((h\otimes k)M)=(h\otimes k)N(M)$ for unitaries $h$ and $k$,
              and $W$ is invariant under these left multiplications;
        \item for $M\in W$, $N(M)=A\otimes B$ with $A\neq0$ and $B\neq0$;
        \item for unitaries $x$ and $y$, $x\otimes y\in W$ and
              $N(x\otimes y)=x\otimes y$.
    \end{enumerate}
    This is the second step of the direct construction that replaces the
    slice theorem in the proof of \cite[Corollary~IV.7]{Cirac2017MPU}. The gap
    is the closed interval $[a,b]$, since avoiding an open interval is not an
    open condition on the spectrum.
    % Source lines: paper_v2.tex 942--947.
\end{lemma}

\begin{proof}\leanok
    \uses{lem:matrix_realignment_kronecker}
    The map $P$ is continuous by continuity of the functional calculus, and
    $W$ is open because the spectrum is upper semicontinuous. Left
    multiplication by $h\otimes k$ sends $R(M)$ to
    $(h\otimes\mathbb 1)R(M)(k\otimes\mathbb 1)^{\mathsf T}$, so it conjugates
    $R(M)R(M)^\dagger$ and $P(M)$ by the unitary $h\otimes\mathbb 1$. On $W$ the
    function $f$ takes only the values $0$ and $1$ on the spectrum, so $P(M)$ is
    an orthogonal projector whose rank equals its trace, which is $1$. Hence
    $P(M)R(M)$ has rank at most one; it is nonzero, for otherwise
    $f(t)t=0$ on the spectrum, which forces $P(M)=0$. The realignment lemma
    then writes $N(M)$ as a Kronecker product of nonzero factors. Finally,
    $R(\mathbb 1)R(\mathbb 1)^\dagger$ is $\ell r$ times a rank-one projector,
    so $P(\mathbb 1)R(\mathbb 1)=R(\mathbb 1)$ and
    $N(\mathbb 1)=\mathbb 1$; equivariance gives the statement for
    $x\otimes y$.
\end{proof}

\begin{lemma}[A local section for local unitaries]
    \label{lem:kronecker_local_section}
    \lean{Matrix.unitaryPart, Matrix.continuousOn_unitaryPart,
        Matrix.unitaryPart_mem_unitaryGroup, Matrix.unitaryPart_unitary_mul,
        Matrix.unitaryPart_kronecker, Matrix.localSection,
        Matrix.localSectionDomain, Matrix.isOpen_localSectionDomain,
        Matrix.continuousOn_localSection, Matrix.localSection_kronecker_mul,
        Matrix.exists_localSection_eq_kronecker_mul,
        Matrix.kronecker_mul_mem_localSectionDomain,
        Matrix.localSection_kronecker_mul_self}
    \leanok
    \uses{lem:nearest_kronecker_product}
    For an invertible matrix $X$ let $Q(X)=X(X^\dagger X)^{-1/2}$. Then $Q$ is
    continuous on the invertible matrices and unitary-valued, $Q(gX)=gQ(X)$ for
    unitary $g$, and $Q(A\otimes B)=Q(A)\otimes Q(B)$ for invertible $A$, $B$.
    Let $0<a<b<\ell r$, let $N$ and the gap set be as in
    Lemma~\ref{lem:nearest_kronecker_product}, and fix a unitary $u_0$. Put
    $g(u)=Q(N(uu_0^\dagger))$ and $\sigma(u)=g(u)^\dagger u$ on the set $W$ of
    those $u$ for which $uu_0^\dagger$ lies in the gap set and
    $N(uu_0^\dagger)$ is invertible. Then $W$ is open, invariant under left
    multiplication by $x\otimes y$ for unitaries $x$, $y$, and contains every
    $(x\otimes y)u_0$; the map $\sigma$ is continuous on $W$, satisfies
    $\sigma((x\otimes y)u)=\sigma(u)$, sends each $u\in W$ to $(x\otimes y)u$
    for some unitaries $x$, $y$, and fixes $u_0$ on its orbit:
    $\sigma((x\otimes y)u_0)=u_0$. This is the third step of the direct
    construction that replaces the slice theorem in the proof of
    \cite[Corollary~IV.7]{Cirac2017MPU}.
    % Source lines: paper_v2.tex 942--947.
\end{lemma}

\begin{proof}\leanok
    \uses{lem:nearest_kronecker_product}
    The factor $(X^\dagger X)^{-1/2}$ is a continuous function of the positive
    definite matrix $X^\dagger X$, and $Q(X)^\dagger Q(X)=1$. Since
    $(gX)^\dagger(gX)=X^\dagger X$, $Q$ commutes with unitaries on the left. For
    a Kronecker product,
    $(A\otimes B)^\dagger(A\otimes B)=A^\dagger A\otimes B^\dagger B$, and the
    inverse square root of a Kronecker product of positive definite matrices
    is the Kronecker product of the inverse square roots, by uniqueness of
    positive square roots. On $W$ the nearest Kronecker product writes
    $N(uu_0^\dagger)=A\otimes B$ with $A$, $B$ invertible, so
    $g(u)=Q(A)\otimes Q(B)$ is a product of local unitaries. Equivariance of $N$
    and of $Q$ gives $g((x\otimes y)u)=(x\otimes y)g(u)$, hence the invariance
    of $\sigma$. Finally $N(\mathbb 1)=\mathbb 1$ and $Q(\mathbb 1)=\mathbb 1$
    give $\sigma(u_0)=u_0$.
\end{proof}

\begin{lemma}[Lifting paths continuous modulo local unitaries]
    \label{lem:kronecker_path_lifting}
    \lean{exists_continuousOn_Icc_lift,
        Matrix.exists_continuous_lift_of_continuous_mod_localUnitary}
    \leanok
    \uses{lem:kronecker_local_section}
    Let $u:[0,1]\to U(\ell r)$ be continuous modulo local unitaries: for every
    $p$ and $\varepsilon>0$ there is $\delta>0$ such that for $|q-p|<\delta$
    some unitaries $x$, $y$ give $\|u(q)-(x\otimes y)u(p)\|<\varepsilon$. Then
    there is a continuous $\hat u:[0,1]\to U(\ell r)$ with
    $\hat u(p)=(x\otimes y)u(p)$ for unitaries $x$, $y$ depending on $p$.
    Nothing is assumed about the regularity of $u$. This is the step of the
    proof of \cite[Corollary~IV.7]{Cirac2017MPU} that the source obtains from
    the slice theorem of Montgomery and Yang.
    % Source lines: paper_v2.tex 942--947.
\end{lemma}

\begin{proof}\leanok
    \uses{lem:kronecker_local_section}
    Fix $p$ and let $\sigma_p$ be the local section of
    Lemma~\ref{lem:kronecker_local_section} with base point $u(p)$, defined on
    the open set $W_p$. If $u(q_0)\in W_p$ and $q$ is close to $q_0$, then
    $(x\otimes y)^\dagger u(q)$ is close to $u(q_0)$ for suitable unitaries
    $x$, $y$, so it lies in $W_p$; as $W_p$ and $\sigma_p$ are invariant under
    local unitaries, $u(q)\in W_p$ and
    $\sigma_p(u(q))=\sigma_p((x\otimes y)^\dagger u(q))$ is close to
    $\sigma_p(u(q_0))$. Hence $q\mapsto\sigma_p(u(q))$ is a continuous lift on
    an open neighbourhood of $p$. A Lebesgue number for this cover of $[0,1]$
    divides $[0,1]$ into finitely many intervals, each carrying a continuous
    lift. Proceeding from left to right, the lift on the next interval is
    multiplied by the constant local unitary that matches it to the lift
    already built at the common endpoint; the glued map is continuous and
    remains a lift.
\end{proof}

\begin{lemma}[The functional calculus under unitary conjugation]
    \label{lem:cfc_unitary_conjugation}
    \lean{Matrix.IsHermitian.cfc_eq_of_eq_unitary_conj,
        Matrix.PosSemidef.supportInvSqrt_eq_of_eq_unitary_conj}
    \leanok
    Let $x$ be unitary and let $A$ and $B=xAx^\dagger$ be Hermitian. For every
    real function $f$, $f(B)=xf(A)x^\dagger$. In particular, if $A$ and $B$ are
    positive semidefinite, the inverse square roots on their supports satisfy
    $B^{-1/2}=xA^{-1/2}x^\dagger$. This compares the normalization matrices of
    two compact decompositions of one matrix in
    \cite[lines~479--502]{Cirac2017MPU}.
\end{lemma}

\begin{proof}\leanok
    Conjugation by $x$ is a continuous star-algebra automorphism, and such
    automorphisms commute with the continuous functional calculus.
\end{proof}
